\section{Benchmark}


% 数据是如何构建的
\subsection{Construction}

The construction of our benchmark involves a systematic pipeline that traces model derivations from base models to their descendants on Hugging Face. This process is divided into five key steps:

\paragraph{Step 1: Obtaining Base Model List}
We begin by collecting a list of popular base models for the \texttt{text-generation} task from Hugging Face. Utilizing the \texttt{HfApi} provided by \texttt{huggingface\_hub}, we retrieve the top 50 non-gated models sorted by the number of likes. The model identifiers are stored in a file named \texttt{base\_models.txt}, where each line corresponds to one base model ID. This step is implemented via \textbf{Script 1}.

\paragraph{Step 2: Building the Model Derivation Tree}
Based on \texttt{base\_models.txt}, we recursively crawl derivative models to construct a “model forest”. For each base model, we explore its derivatives up to four levels deep, covering four derivation types: \texttt{finetune}, \texttt{adapter}, \texttt{quantized}, and \texttt{merge}. For each derivation type, a maximum of six sibling nodes are considered. Every node in the tree records metadata such as model name, crawl timestamp, depth level, and its child nodes. The resulting forest structure is serialized into \texttt{model\_forest.json}. This step is implemented via \textbf{Script 2}.

\paragraph{Step 3: Constructing Model Pairs}
From the model forest, we extract two categories of model pairs: adjacent parent-child pairs and non-adjacent ancestor-descendant pairs. By performing a Breadth-First Search (BFS) traversal on each tree, we generate:
\begin{itemize}
    \item \textbf{Adjacent Pairs}: Direct parent-child model pairs.
    \item \textbf{Non-Adjacent Pairs}: All ancestor-descendant pairs where the ancestor is not the immediate parent.
\end{itemize}
These pairs are recorded in \texttt{adjacent\_pairs.jsonl} and \texttt{non\_adjacent\_pairs.jsonl}, respectively. This step is implemented via \textbf{Script 3}.

\paragraph{Step 4: Single-Threaded Model Validation}
To ensure data quality, we validate the structural correctness of all adjacent model pairs. For each model, we inspect its file tree on Hugging Face, checking for the presence of critical files such as \texttt{config.json}, \texttt{adapter\_config.json}, and various parameter files (e.g., \texttt{.bin}, \texttt{.pt}, \texttt{.safetensors}, etc.). We verify whether the claimed derivation type aligns with the actual file structure and whether the parent model is valid. The results are compiled into a comprehensive validation report named \texttt{model\_validation\_report.json}. This step is implemented via \textbf{Script 4}.

\paragraph{Step 5: Filtering Erroneous Model Pairs}
Based on the validation report, we remove all erroneous model pairs from the dataset. Models flagged with any inconsistencies or missing critical files are excluded. The cleaned set of adjacent model pairs is saved in \texttt{adjacent\_pairs\_without\_error.json}. This step is implemented via \textbf{Script 5}.

\subsection{Characteristics}

Our benchmark exhibits several notable characteristics that reflect its diversity and representativeness of the Hugging Face model ecosystem:

\begin{itemize}
    \item \textbf{Variety of Derivation Types}: The dataset covers four major derivation types—\texttt{finetune}, \texttt{adapter}, \texttt{quantized}, and \texttt{merge}—capturing a wide range of model transformation patterns.
    
    \item \textbf{Hierarchical Depth}: With a recursive crawling depth of up to four levels, the benchmark contains both shallow and deep derivation chains, enabling analysis of both immediate and long-range model inheritance.
    
    \item \textbf{Scale and Coverage}: Starting from 50 popular base models, the recursive traversal results in a substantial model forest, ensuring coverage of highly influential models as well as their derivative variants.
    
    \item \textbf{Data Cleanliness}: Through rigorous validation and error filtering, the dataset ensures high structural integrity, making it suitable for tasks that require reliable model lineage information.
    
    \item \textbf{Real-World Relevance}: The benchmark reflects actual user-driven derivation patterns on Hugging Face, providing insights into community practices such as model fine-tuning trends, adapter usage, and quantization adoption.
\end{itemize}

% show the statistics of the benchmark
\begin{table}[t]
    \centering
    \caption{Statistics of the \textsc{GrandMoTHer} benchmark.}
    \label{tab:grandmother-statistics}
    \small
    \begin{tabular}{l r}
        \toprule
        \textbf{Statistic} & \textbf{Value} \\
        \midrule
        \#Graphs & 31 \\
        \#Model pairs (edges) & 248 \\
        Min / Median / Max pairs per graph & 1 / 6 / 22 \\
        Avg.\ pairs per graph & 8.0 \\
        \#Models (nodes) & 279 \\
        \#Root models & 31 \\
        \#Intermediate models & 44 \\
        \#Leaf models & 204 \\
        Share of leaf models & 73.1\% \\
        Min / Avg / Max models per graph & 2 / 9.0 / 23 \\
        Avg.\ out-degree of non-leaf models & 3.31 \\
        Max fine-tuning depth & 2 \\
        \bottomrule
    \end{tabular}
\end{table}